% Round 16. Reconstructed framed scalar comparison, independently reviewed.
% This is a comparison on supplied regular numerically admissible model data.
\section{Comparison with a supplied regular numerically admissible model}
\label{r16mm:section}

The complete mathematical source passed independent trace-branch
review at SHA256
\texttt{0dc8be905b0bf30b1dbf15eaaa75acdec5e9950045dbe3d41fe28b108de0df38},
recorded in
\path{higher_trace/regular_admissible_model_comparison_independent_audit.md}.
The current text adds the review's optional explicit cap/Adams
justification and this status notice to that core proof. The final
abelian identification corollary was independently checked by the
model branch at SHA256
\texttt{89c45cb47dc6c74b73f5b0361d8ba52bcae638ff36f3e488c2cc010351df8a12}.
It identifies an already published classical height; it is not
a new existence claim for the classical abelian pairing.

The historical Round 7 model-comparison manuscripts retained an
unresolved full finite-trace interface. The present argument rebuilds
the fixed-supported-complex scalar comparison using the independently
proved R11 construction and H16.1. It then checks the algebraic
comparison with the finite Gillet--Soul\'e term directly. Historical
proof statuses are not changed. No existence of an admissible model
lift for a general class is inferred.

Let $K/\mathbf Q_p$ have valuation ring $R$, residue field $k$,
uniformizer $\varpi$, and prescribed scalar
$c_R=-\log|\varpi|>0$. All analytic integrations use the native
$K$-weights. Every degree of a zero-cycle below is taken on a
proper vertically supported scheme over $k$, with nilpotent
thickenings identified with their reductions for Chow groups.

\begin{lemma}[Native vertical Cartier evaluation]
\label{r16mm:vertical}
Let $\mathcal V/R$ be integral, normal, flat and projective,
of relative dimension $t$, with normal generic fiber $V$.
Write $[\mathcal V_k]=\sum_i b_i[Y_i]$. Let $E$ be a vertical
Cartier divisor and write its Weil cycle as $\sum_i a_i[Y_i]$.
For a degree-$t$ polynomial $P$ in first Chern classes of model
line bundles, let $P^{\rm an}$ be its actual model form and
$g_E=-\log\|1\|_{\mathcal O(E)}$. Then
\[
 \int_{V^{\rm an}}g_E\,P^{\rm an}
     =c_R\sum_i a_i\deg_k(P|_{Y_i}\cap[Y_i])
     =c_R\deg_k(P|_E\cap[E]).
 \tag{MM16.1}\label{r16mm:vertical-eq}
\]
For a signed divisor the final expression is interpreted by
linearity on its positive and negative supports.
\end{lemma}

\begin{proof}
The actual-product/Bedford--Taylor comparison in the frozen
R10 model-Chern theorem identifies the top actual measure with
the native model measure. Chambert-Loir--Ducros, the manuscript
dated 25 July 2025, Theorem~18.5.5 and Proposition~18.5.8, give
\[
 [P^{\rm an}]=\sum_i b_i\deg_k(P|_{Y_i}\cap[Y_i])\,\delta_{\xi_i},
\]
where $\xi_i$ is the unique divisorial analytic point over the
generic point of $Y_i$. Normality and excellence allow the
relatively normal formal completion used in those statements.
The generic multiplicity and defect are one, and their valuation
index is exactly $b_i$, not one. The same frozen comparison works
for the line bundles pulled to this normal model: locally
$[uT]=u[T]$ and $[Du\wedge T]=D(u[T])$ for the continuous model
weight $u$ and a closed model product $T$. This is the precise
Bedford--Taylor recursion, so the actual and classical model
measures coincide without changing the field.

At $\xi_i$, the prescribed valuation is $c_R/b_i$ times the
integer divisor valuation. Thus $g_E(\xi_i)=c_Ra_i/b_i$.
Multiplication by the preceding finite signed measure cancels
$b_i$ and proves \eqref{r16mm:vertical-eq}. Continuous scalar
multiplication is permitted for that Radon measure; the scalar
here is a global actual PL function. All residue-field degrees
remain inside $\deg_k$.
\end{proof}

\begin{lemma}[The supported line-complex character over the trait]
\label{r16mm:cartier}
Let $\mathcal V$ be integral projective over $R$, of total
dimension $a$, and let $s:L\to M$ be a nonzero line-bundle
map with Cartier zero scheme $E$. With $M$ in homological
degree zero, for every $r\geq1$ one has
\[
 \operatorname{ch}_r^E([L\xrightarrow{s}M])\cap[\mathcal V]
  =\frac1{r!}\sum_{j=0}^{r-1}
        c_1(M|_E)^{r-1-j}c_1(L|_E)^j\cap[E].
 \tag{MM16.2}\label{r16mm:cartier-eq}
\]
This is a supported Chow identity over the excellent DVR itself,
including entirely vertical and nonreduced Cartier schemes.
\end{lemma}

\begin{proof}
Put $N=M\otimes L^{-1}=\mathcal O(E)$. Blow up
$\mathbf P^1_{\mathcal V}$ along $E\times\{\infty\}$.
Its infinity fiber consists of the strict transform
$\mathcal V'$ and $P=\mathbf P_E(N^\vee|_E\oplus\mathcal O_E)$.
The strict transform $J$ of $\mathbf P^1_E$ is Cartier, disjoint
from $\mathcal V'$, and meets $P$ in the zero section of the
normal line $N|_E$. Locally the relevant chart is
$A[t,y]/(ty-f)$ with $J=(y)$, where $f$ is a nonzerodivisor;
these assertions therefore also hold when the ambient model
or $E$ is singular.

The perfect complex $q^*M\otimes[\mathcal O(-J)\to\mathcal O]$
pulls back the original complex away from infinity and is zero
on $\mathcal V'$. The general-base deformation lemma Stacks
0FE8 (42.51.2), with the geometry of 0FBG, computes its localized
character by the exceptional contribution. On $P$, put
$m=c_1(M|_E)$, $l=c_1(L|_E)$, $\nu=m-l$, and let $h$ be
the zero-section divisor class. Its normal bundle is $N|_E$,
so $h^2=\nu h$, $\pi_*(h\cap[P])=[E]$, and $\pi_*[P]=0$.
Therefore the exceptional push is the formal polynomial series
\[
       \pi_*\bigl(e^m(1-e^{-h})\cap[P]\bigr)
                   =\frac{e^m-e^l}{m-l}\cap[E].
\]
Its degree $r-1$ term is \eqref{r16mm:cartier-eq}. The quotient
denotes its polynomial series, never inversion of a Chow class.
The push is to $E$, so the support has not been forgotten.
This reconstructs the general-base Cartier step of the R7
manuscript without importing a field-only Riemann--Roch theorem.
\end{proof}

\begin{theorem}[Framed scalar comparison for the concrete R11 Green pair]
\label{r16mm:framed}
Let $\mathcal X/R$ be integral flat projective with smooth
geometrically integral generic fiber $X$. Let $C$ be a bounded
model vector-bundle complex exact off a model closed support
$\mathcal Z$. Fix $1\leq r\leq\dim X$, and assume its generic
support $Z$ is pure of codimension $r$ or is empty. Let $G_C$
be its degree-$(r-1,r-1)$ actual form from the
concrete R11 line-complex construction. For an integral
$(r-1)$-dimensional subvariety $W\subset X\setminus Z$, let
$\mathcal W$ be its schematic closure and
$B=\mathcal Z\cap\mathcal W$. Then
\[
 \int_{W^{\rm an}}G_C
       =c_R\deg_k\left(
          \operatorname{ch}_r^B(C|^{\mathbf L}_{\mathcal W})
                                  \cap[\mathcal W]\right).
 \tag{MM16.3}\label{r16mm:framed-eq}
\]
For empty $Z$, use the frozen generic-acyclic version of that
construction, so $G_C$ is global actual. The comparison is for
the fixed supported perfect complex, with
its differential and vertical support retained. It includes a
testing cycle contained in the branch locus of the chosen finite
presentation and arbitrary vertical excess on the model.
\end{theorem}

\begin{proof}
Use the R11 supported line-complex presentation on a proper
model map $f:\mathcal T\to\mathcal X$, whose generic fiber
has degree $d$ and is finite flat over $X\setminus Z$.
It gives, in the supported perfect group,
$[f^*C]=\sum_\alpha\epsilon_\alpha[L_\alpha\to M_\alpha]$
and, over that complement,
\[
 G_C=d^{-1}\operatorname{Tr}_f G_T,
 \qquad
 G_T=\sum_\alpha\epsilon_\alpha g_\alpha
                   B_{r-1}(l_\alpha,m_\alpha),
 \quad
 B_{r-1}(l,m)=\frac1{r!}\sum_{j=0}^{r-1}m^{r-1-j}l^j.
 \tag{MM16.4}\label{r16mm:presentation}
\]
Retain the full finite fiber over the normalization of $W$ and
its generic Artinian lengths $e_\lambda$. Let $V_\lambda$ be
the resulting normal generic components, and let
$\mathcal V_\lambda$ be their normal projective closures on
the presentation model, using a graph closure if needed. They
are flat over $R$ and proper over $\mathcal W$. Their weighted
fundamental cycles satisfy
\[
 \sum_\lambda e_\lambda
                   (h_\lambda)_*[\mathcal V_\lambda]
                                      =d[\mathcal W].
 \tag{MM16.5}\label{r16mm:closure-degree}
\]
This is an equality of model cycles because the closures are
integral and dominate $\mathcal W$ with those generic degrees.
No vertical component of a possibly nonflat full model fiber
product has been inserted or silently deleted.

H16.1 gives the actual finite-flat base-change formula on the
normalization of $W$, including all $e_\lambda$. Integrating
its actual top form and using finite adjunction gives
\[
 \int_WG_C=d^{-1}\sum_{\lambda,\alpha}
 e_\lambda\epsilon_\alpha
       \int_{V_\lambda}g_{\alpha\lambda}
                          B_{r-1}(l_{\alpha\lambda},m_{\alpha\lambda}).
 \tag{MM16.6}\label{r16mm:trace-scalar}
\]
Top actual integration is unchanged under normalization.
Compactness is automatic for these proper generic cycles;
equivalently multiply by an ordinary smooth compact scalar
equal to one near $W$ inside $X\setminus Z$ before applying
the finite adjunction. No individual geometric sheet is used.

Each pulled line map is an isomorphism on $V_\lambda$, so its
Cartier divisor $E_{\alpha\lambda}$ on $\mathcal V_\lambda$
is entirely vertical. Lemma~\ref{r16mm:vertical} computes each
integral in \eqref{r16mm:trace-scalar} as $c_R$ times the
supported degree of $B_{r-1}\cap[E_{\alpha\lambda}]$.
Lemma~\ref{r16mm:cartier} identifies it with the localized
degree-$r$ character of that line complex. Additivity and
arbitrary derived base change of localized characters identify
the sum over $\alpha$ with
\[
 \operatorname{ch}_r^{h_\lambda^{-1}B}
       (h_\lambda^*(C|^{\mathbf L}_{\mathcal W}))
                                         \cap[\mathcal V_\lambda].
\]
The supported line-complex identity pulls back because its
termwise vector-bundle sequences are locally split. A common
nilpotent support thickening can be retained throughout.
Proper bivariant projection and \eqref{r16mm:closure-degree}
then make the weighted pushed sum exactly
$d\operatorname{ch}_r^B(C|^{\mathbf L}_{\mathcal W})\cap[\mathcal W]$.
All these pushes in the degree calculation are between proper
vertical supports. Divide by $d$ to obtain the formula.
\end{proof}

\begin{lemma}[Supported character degree and finite Euler intersection]
\label{r16mm:euler}
Let $\mathcal X/R$ be regular integral flat projective of total
dimension $N=n+1$. Let $\mathcal Z,\mathcal W$ be integral
model cycles of codimensions $r,s$ with $r+s=N$, whose generic
fibers are disjoint. Put $B=\mathcal Z\cap\mathcal W$ and
choose finite vector-bundle resolutions $C_Z,C_W$. Then
\[
 \deg_k\left(\operatorname{ch}_r^B
          (C_Z|^{\mathbf L}_{\mathcal W})\cap[\mathcal W]\right)
 =\chi_R R\Gamma\left(\mathcal X,
         \mathcal O_{\mathcal Z}\otimes^{\mathbf L}
                                   \mathcal O_{\mathcal W}\right),
 \tag{MM16.7}\label{r16mm:euler-eq}
\]
where $\chi_R$ is the alternating sum of finite $R$-lengths.
The intersection $B$ may have positive dimension.
\end{lemma}

\begin{proof}
This is the supported calculation in the R7 finite-bridge
manuscript; its primary algebraic inputs can be checked directly.
Write $a_i=\operatorname{ch}_i^{\mathcal Z}(C_Z)$ and
$b_j=\operatorname{ch}_j^{\mathcal W}(C_W)$.
Support dimension gives $a_i\cap[\mathcal X]=0$ for $i<r$
in the Chow group of $\mathcal Z$ itself, and similarly for
$j<s$. The leading terms are $[\mathcal Z]$ and
$[\mathcal W]$ by Stacks 0FEW. Centrality (42.52.2) permits
each lower term to act first, so its product vanishes before
forgetting intersection support. Intersected-support
multiplicativity (42.52.6) consequently gives
\[
 \operatorname{ch}^B(C_Z\otimes C_W)\cap[\mathcal X]
  =a_r\cap[\mathcal W]\in\operatorname{CH}_0(B)_{\mathbf Q},
 \tag{MM16.8}\label{r16mm:pure-zero}
\]
with no other-dimensional terms: total degree below $N$ has
a lower factor, and total degree above $N$ has negative output
dimension.

For the closed immersion $B\to\mathcal X$, Stacks 0FEW gives
the isomorphism
\[
 G_0(B)_{\mathbf Q}\xrightarrow{\sim}\operatorname{CH}_*(B)_{\mathbf Q},
 \qquad [F]\longmapsto
           \operatorname{ch}^B(F)\cap[\mathcal X].
\]
It is stated over the general excellent base, with regular
quasi-compact ambient scheme. It sends the class of a closed
point sheaf to that point, since all lower-dimensional terms
would have negative dimension. If a zero-cycle represents the
right side of \eqref{r16mm:pure-zero}, injectivity therefore
identifies the alternating cohomology-sheaf class of
$C_Z\otimes C_W$ with the corresponding point-sheaf combination.
The proper Euler-length map is additive on $G_0(B)$ and takes
$\mathcal O_x$ to $[\kappa(x):k]$. Applying it proves
\eqref{r16mm:euler-eq}. Coherent cohomology has finite length
because $B$ is vertical and proper. No Todd correction is
discarded: the supported character class has already been
proved purely zero-dimensional. No degree on the Chow group of
the whole trait model is used.
\end{proof}

\begin{theorem}[Local comparison on supplied numerical model lifts]
\label{r16mm:local}
Let $\mathcal X/R$ be regular integral flat projective, with
smooth geometrically integral generic fiber $X$ of dimension
$n$. Let $Z,W$ be disjoint geometrically homologically trivial
rational cycles of codimensions $r,s$, $r+s=n+1$.
Suppose rational model cycles $\Theta_Z,\Theta_W$ of those
codimensions have generic cycles exactly $Z,W$ and are
numerically admissible: their operational intersection degree
against every vertical cycle of the complementary testing
dimension is zero on each special component. Then both cycles
have exact strong ordinary actual Green forms, and their
choice-independent local pairing satisfies
\[
 \lambda_K(Z,W)=c_R\,I_R(\Theta_Z,\Theta_W),
 \tag{MM16.9}\label{r16mm:local-eq}
\]
where $I_R$ is the bilinear extension of the finite Euler tensor
intersection in \eqref{r16mm:euler-eq}. The statement permits
singular cycle models, nonreduced special fibers and arbitrary
vertical excess. It makes no assertion that such lifts always
exist.
\end{theorem}

\begin{proof}
Write each specified model cycle as a rational sum of prime
cycles and choose vector-bundle resolutions of their structure
sheaves; regularity and projectivity supply such resolutions.
Their degree-$r$ Chern characters represent the operational
class of $\Theta_Z$: the localized leading-term identity gives
this equality after capping with the regular model, and the frozen
R12 cap/Adams characterization identifies that degree-$r$ term
with its operational class. The strict R10 model-Chern identities and
numerical detector make the sum of their actual curvatures
literally zero. Sum their concrete R11 Green forms. A vertical
component contributes a global actual form from the R11
generic-acyclic construction. The horizontal components are
exactly the closures of the final nonzero prime coefficients
of $Z$. Hence the sum is actual on the exact complement of
$Z$ and has its strong residue on all actual tests. The same
argument applies to $W$.

By Theorem~\ref{r16mm:framed} and Lemma~\ref{r16mm:euler},
the value of this constructed form on $W$ is
$c_R I_R(\Theta_Z,\overline W)$.
Write $\Theta_W=\overline W+V$, where $V$ is a rational
vertical cycle of dimension $r$. The additional localized
degree against each prime component of $V$ pushes into that
proper $k$-variety as the ordinary operational degree of
$\Theta_Z$ on it. Numerical admissibility makes this degree
zero. Thus $\overline W$ may be replaced by $\Theta_W$.
Disjoint-polar choice independence, already proved because both
generic cycles have strong G, replaces the constructed form
by any supplied exact strong Green form. This proves
\eqref{r16mm:local-eq} without selecting an unspecified global
actual primitive whose scalar value might be uncontrolled.
\end{proof}

\begin{theorem}[Global comparison with the classical admissible model height]
\label{r16mm:global}
Let $\mathcal X/\mathcal O_F$ be regular integral flat projective
with smooth projective geometrically integral generic fiber $X$.
Let $z,w$ be complementary geometrically homologically trivial
rational Chow classes admitting numerically admissible rational
model Chow lifts $\Theta,\Phi$ on this supplied model. Then
they belong to the domains $\mathcal H$ of
Theorem~\ref{r16h:global-height}, and
\[
 h_X^{\rm an}(z,w)=\frac1{[F:\mathbf Q]}
                \widehat{\deg}\bigl(\widehat\Theta\cdot\widehat\Phi\bigr),
 \tag{MM16.10}\label{r16mm:global-eq}
\]
where the arithmetic lifts have zero archimedean curvature.
The right side is the classical admissible intersection height
on the supplied domain. The comparison is for the fixed
absolute-value normalization of the analytic height.
\end{theorem}

\begin{proof}
Choose disjoint generic cycle representatives $Z,W$ of $z,w$.
Each given model Chow lift has a rational cycle representative
whose generic cycle is exactly the chosen one: lift the finitely
many rational-equivalence functions from their generic carriers
to their model closures and retain all their vertical divisor
terms. This changes only the model representative, not its
Chow class or its numerical admissibility. There is no requirement
to move the resulting special-fiber intersections into proper
position.

After completion at a finite place $v$, regularity is preserved
and the same numerical and supported identities apply. Local
numerical lifts also imply obstruction zero by the frozen
R11 criterion, hence membership in $\mathcal H$. The preceding
theorem identifies the analytic local value with $c_v I_v$.
If $v\mid p$ has ramification index $e_v$ and residue degree
$f_v$, then $c_v=(\log p)/e_v$ and
\[
 \frac{[F_v:\mathbf Q_p]}{[F:\mathbf Q]}c_v I_v
             =\frac{\log\#k_v}{[F:\mathbf Q]}I_v.
 \tag{MM16.11}\label{r16mm:weight}
\]
The right side is exactly the finite arithmetic-degree term:
Gillet--Soul\'e, Theorem~4.3.2 and Remark~4.3.8(iv), define it
by supported derived tensor, proper cohomology and length.
The explicit double sum of lengths is
$\sum_{a,b}(-1)^{a+b}\operatorname{length}_{R_v}
H^a(\mathcal X_{R_v},\mathcal{T}or_b(\mathcal O_{\Theta},
\mathcal O_{\Phi}))$, extended bilinearly in prime model cycles.

At infinity choose logarithmic-type zero-curvature complex
Green forms. In Gillet--Soul\'e's principal normalization
$g_Z=-\log|f|^2$, the corresponding form for our normalization
is $g_Z/2$. Since the second Green curvature is zero, the
archimedean star-product integral is its evaluation on the
other cycle. The arithmetic degree contributes
$\frac1{2[F:\mathbf Q]}\sum_{\sigma:F\hookrightarrow\mathbf C}
\int_{W_\sigma}g_{Z,\sigma}$, which is precisely our sum over
archimedean places with weights one at a real place and two
at a complex place, divided by $[F:\mathbf Q]$.

These are the finite and infinite terms of
\eqref{r16mm:global-eq}. Only finitely many finite terms occur,
either by proper vertically supported cohomology of the
generic-disjoint pair on the arithmetic model, or by the
already proved height finiteness theorem. The latter theorem
also removes all generic moving and Green choices. The usual
arithmetic Chow identity for the product with a smooth form
class shows that zero-curvature arithmetic lifts have the
same degree product. Thus no unrecorded Bott--Chern correction
or rational-equivalence logarithm is omitted.
\end{proof}

\begin{corollary}[Identification with the published abelian height]
\label{r16mm:abelian}
Let $A/F$ be an abelian variety over a number field. In every
complementary pair of codimensions, the full analytic height
supplied by \texttt{r16ab:local-G} and
Theorem~\ref{r16h:global-height} equals K\"unnemann's classical
abelian height divided by $[F:\mathbf Q]$, with the same
Gillet--Soul\'e product and scalar sign convention.
\end{corollary}

\begin{proof}
After a finite extension $E/F$, semi-abelian reduction holds.
K\"unnemann, \emph{Height pairings for algebraic cycles on abelian
varieties}, Theorem~1.6, supplies a regular flat projective
model over $\mathcal O_E$ and lifts of all geometric homological
classes orthogonal to all vertical classes. A vertical testing
cycle supported at a single finite place contributes only its
local intersection degree times the positive factor
$\log\#k_v$. Thus this orthogonality is precisely the
placewise numerical admissibility needed in
Theorem~\ref{r16mm:global}. Every cycle on a completed special
fiber is a cycle on that same algebraic special fiber, so the
condition includes all its component tests.

Theorem~\ref{r16mm:global} identifies the analytic pairing over
$E$ with that model pairing divided by $[E:\mathbf Q]$.
K\"unnemann's Section~1.4, equation~(7), and Lemma~1.5 descend
the classical unnormalized pairing by $1/[E:F]$; its Corollary~1.7
is the resulting published abelian height over $F$.
Theorem~\ref{r16h:global-height} makes the analytic normalized
pairing invariant under this same finite extension. Since
$[E:\mathbf Q]=[E:F][F:\mathbf Q]$, the two formulas give the
assertion. Only a constant-field extension and its native degree
are used, not an arithmetic alteration comparison.
\end{proof}

\paragraph{The framing limitation and the remaining comparison boundary.}
For $C=[\mathcal O\xrightarrow{\varpi^a}\mathcal O]$ in degree
one, the absolute model $K_0$ class and actual curvature are
zero, but the constructed scalar is $a c_R$. Its supported
character has exactly degree $a$ on a degree-one point model.
Thus curvature equality alone cannot prove a model scalar
comparison on arbitrary test cycles. Theorem~\ref{r16mm:framed}
keeps this framing; Theorem~\ref{r16mm:local} removes it only
through the already proved two-sided Green-choice independence.

The global theorem compares with a supplied regular numerical
admissible model. It does not assert the existence of such a
global model or admissible lifts for every class in $\mathcal H$,
and it does not promote a bare strong-G solution to obstruction
zero. No arithmetic-alteration projection formula is folded into
\eqref{r16mm:global-eq}. A separate scalar descent using the new
Chow-correspondence adjunction is stated in
\texttt{r16ha:alteration-comparison}; that statement still requires
the regular admissible lifts upstairs and supplies its own generic
degree calculation. Parameter-
space continuity of a biextension metric, positivity and
nondegeneracy are not conclusions of this scalar theorem.

\paragraph{Sources and provenance.}
The new analytic substitutions are H16.1 and the final R11
\path{perfect_complex/supported_complex_green.tex}; model
identities are the independently audited R10
\path{chern/model_chern_without_trace.tex}.
The historical arguments reconstructed here are
\path{history_01_07/bb_round7/proofs/good_relative_complex_local_scalar_M.tex}
and
\path{history_01_07/bb_round7/proofs/good_localized_chern_GS_finite_bridge.tex}.
Their prior review flags are not removed.
Primary sources re-opened for this reconstruction are
\url{https://www.imj-prg.fr/~antoine.ducros/fdceb.pdf},
Theorem~18.5.5 and Proposition~18.5.8;
\url{https://stacks.math.columbia.edu/tag/0FE8},
\url{https://stacks.math.columbia.edu/tag/0FB9},
\url{https://stacks.math.columbia.edu/tag/0FEW}; and
Gillet--Soul\'e,
\url{https://www.numdam.org/article/PMIHES_1990__72__93_0.pdf},
Theorem~4.3.2 and Remark~4.3.8(iv)--(v).
